% test-013.tex -- comandos \sprecta, \sprayo, \spside y \spmside
%
% Compilar:  lualatex-dev test-013.tex
% Validar:   show-pdf-tags --xml test-013.pdf | rnv-wrapp
%            verapdf --flavour ua2 --format text test-013.pdf
%
% Cada caso es un parrafo numerado, para poder referirse a el al escuchar
% el PDF. Los comentarios "doc:" indican lo que documenta el manual; en el
% resto, la lectura se revisa escuchando. Todos trabajan en modo
% matematico. \sprecta recibe una letra o dos puntos; \sprayo, \spside y
% \spmside, dos puntos. Un punto es una letra mayuscula con prima o con
% subindice entre llaves. Las entradas invalidas (que detienen la
% compilacion) no van aqui, entre ellas dos letras minusculas juntas y A_1
% sin llaves.
%
% Casos pensados para escuchar: los de silent y concept de \sprecta con
% letra sola, con subindice y con prima (sprecta-silent.patch), y los de
% prefix de \spside y \spmside.
\DocumentMetadata
  {
    lang=es-CL,
    pdfversion=2.0,
    pdfstandard=ua-2,
    tagging=on,
    tagging-setup={math/setup={mathml-SE}},
  }
\documentclass{article}
\usepackage[spanish]{babel}
\usepackage{unicode-math}
\usepackage{spintent}
\usepackage{hyperref}
\hypersetup
  {
    pdfdisplaydoctitle = true,
    pdftitle           = {test-013: sprecta, sprayo, spside y spmside},
  }
\pagestyle{empty}
\begin{document}

\section*{A. Comando sprecta: argumentos}

% doc: «recta A B»
1. Dos puntos juntos: $\sprecta{AB}$.

% doc: «recta m»
2. Una letra minúscula: $\sprecta{m}$.

% doc: «recta A B»
3. Dos puntos con coma y espacio: $\sprecta{A, B}$.

% doc: «recta A B»
4. Dos puntos con coma: $\sprecta{A,B}$.

% doc: «recta A prima B prima»
5. Con primas: $\sprecta{A'B'}$.

% doc: «recta A sub uno B sub dos»
6. Con subíndices: $\sprecta{A_{1}B_{2}}$.

% doc: «recta A sub uno B sub dos»
7. Con subíndices y coma: $\sprecta{A_{1},B_{2}}$.

% doc: «recta A doble prima B»
8. Con dos primas: $\sprecta{A''B}$.

% doc: «recta A sub 12 B sub tres»
9. Con subíndices de dos cifras: $\sprecta{A_{12}B_{3}}$.

% doc: «recta L»
10. Una letra mayúscula: $\sprecta{L}$.

% doc: «recta r sub uno»
11. Una letra con subíndice: $\sprecta{r_{1}}$.

% doc: «recta L sub uno»
12. Una mayúscula con subíndice: $\sprecta{L_{1}}$.

% doc: «recta m prima»
13. Una letra con prima: $\sprecta{m'}$.

% doc: «recta l doble prima»
14. Una letra con dos primas: $\sprecta{l''}$.

\section*{B. Comando sprecta: llaves}

% doc: «recta A B»
15. Con read-cmd school, el valor por defecto: $\sprecta[read-cmd=school]{AB}$.

% doc: «recta a b»
16. Con read-cmd mathcat: $\sprecta[read-cmd=mathcat]{AB}$.

% doc: «recta l»
17. Con read-cmd mathcat y una mayúscula: $\sprecta[read-cmd=mathcat]{L}$.

% doc: «la recta A B»
18. Con prefix y dos puntos: $\sprecta[prefix={la}]{AB}$.

% doc: «la recta m»
19. Con prefix y una letra: $\sprecta[prefix={la}]{m}$.

% doc: «la recta r sub uno»
20. Con prefix y una letra con subíndice: $\sprecta[prefix={la}]{r_{1}}$.

% doc: «la recta m prima»
21. Con prefix y una letra con prima: $\sprecta[prefix={la}]{m'}$.

% doc: «recta A B»
22. Con concept true, el valor por defecto: $\sprecta[concept=true]{AB}$.

% doc: «A B»
23. Con concept false y dos puntos: $\sprecta[concept=false]{AB}$.

% doc: «m»
24. Con concept false y una letra: $\sprecta[concept=false]{m}$.

% doc: «r sub uno»
25. Con concept false y una letra con subíndice: $\sprecta[concept=false]{r_{1}}$.

% doc: «m prima»
26. Con concept false y una letra con prima: $\sprecta[concept=false]{m'}$.

% doc: «recta A B»
27. Con silent false, el valor por defecto: $\sprecta[silent=false]{AB}$.

% doc: «»
28. Con silent true y dos puntos: $\sprecta[silent=true]{AB}$.

% doc: «»
29. Con silent true y una letra: $\sprecta[silent=true]{m}$.

% doc: «»
30. Con silent true y una letra con subíndice: $\sprecta[silent=true]{r_{1}}$.

% doc: «»
31. Con silent true y una letra con prima: $\sprecta[silent=true]{m'}$.

% doc: «»
32. Con silent true y una letra con dos primas: $\sprecta[silent=true]{l''}$.

% doc: «»
33. Con silent true y read-cmd mathcat: $\sprecta[silent=true,read-cmd=mathcat]{r_{1}}$.

% doc: «la recta a subíndice 1 b subíndice 2»
34. Con todas las llaves: $\sprecta[read-cmd=mathcat,prefix={la},concept=true,silent=false]{A_{1}B_{2}}$.

\section*{C. Comando sprayo}

% doc: «rayo A B»
35. Dos puntos juntos: $\sprayo{AB}$.

% doc: «rayo A B»
36. Dos puntos con coma y espacio: $\sprayo{A, B}$.

% doc: «rayo A B»
37. Dos puntos con coma: $\sprayo{A,B}$.

% doc: «rayo A prima B prima»
38. Con primas: $\sprayo{A'B'}$.

% doc: «rayo A sub uno B sub dos»
39. Con subíndices: $\sprayo{A_{1}B_{2}}$.

% doc: «rayo A sub uno B sub dos»
40. Con subíndices y coma: $\sprayo{A_{1},B_{2}}$.

% doc: «rayo O A»
41. Con el origen: $\sprayo{OA}$.

% doc: «rayo A B»
42. Con read-cmd school, el valor por defecto: $\sprayo[read-cmd=school]{AB}$.

% doc: «rayo a b»
43. Con read-cmd mathcat: $\sprayo[read-cmd=mathcat]{AB}$.

% doc: «rayo A B»
44. Con read-sty rayo, el valor por defecto: $\sprayo[read-sty=rayo]{AB}$.

% doc: «semirrecta A B»
45. Con read-sty semirrecta: $\sprayo[read-sty=semirrecta]{AB}$.

% doc: «el rayo A B»
46. Con prefix y read-sty rayo: $\sprayo[prefix={el}]{AB}$.

% doc: «la semirrecta A B»
47. Con prefix y read-sty semirrecta: $\sprayo[prefix={la},read-sty=semirrecta]{AB}$.

% doc: «rayo A B»
48. Con concept true, el valor por defecto: $\sprayo[concept=true]{AB}$.

% doc: «A B»
49. Con concept false: $\sprayo[concept=false]{AB}$.

% doc: «A B»
50. Con concept false y semirrecta: $\sprayo[concept=false,read-sty=semirrecta]{AB}$.

% doc: «rayo A B»
51. Con silent false, el valor por defecto: $\sprayo[silent=false]{AB}$.

% doc: «»
52. Con silent true: $\sprayo[silent=true]{AB}$.

% doc: «»
53. Con silent true y subíndices: $\sprayo[silent=true]{A_{1}B_{2}}$.

% doc: «la semirrecta a subíndice 1 b subíndice 2»
54. Con todas las llaves: $\sprayo[read-cmd=mathcat,read-sty=semirrecta,prefix={la},concept=true,silent=false]{A_{1}B_{2}}$.

\section*{D. Comando spside}

% doc: «lado A B»
55. Dos puntos juntos: $\spside{AB}$.

% doc: «lado A B»
56. Dos puntos con coma y espacio: $\spside{A, B}$.

% doc: «lado A B»
57. Dos puntos con coma: $\spside{A,B}$.

% doc: «lado A prima B prima»
58. Con primas: $\spside{A'B'}$.

% doc: «lado A sub uno B sub dos»
59. Con subíndices: $\spside{A_{1}B_{2}}$.

% doc: «lado A sub uno B sub dos»
60. Con subíndices y coma: $\spside{A_{1},B_{2}}$.

% doc: «lado B C»
61. Con el vértice: $\spside{BC}$.

% doc: «lado A B»
62. Con read-cmd school, el valor por defecto: $\spside[read-cmd=school]{AB}$.

% doc: «lado a b»
63. Con read-cmd mathcat: $\spside[read-cmd=mathcat]{AB}$.

% doc: «lado A B»
64. Con read-sty lado: $\spside[read-sty=lado]{AB}$.

% doc: «segmento A B»
65. Con read-sty segmento: $\spside[read-sty=segmento]{AB}$.

% doc: «trazo A B»
66. Con read-sty trazo: $\spside[read-sty=trazo]{AB}$.

% doc: «el lado A B»
67. Con prefix: $\spside[prefix={el}]{AB}$.

% doc: «el segmento A B»
68. Con prefix y read-sty segmento: $\spside[prefix={el},read-sty=segmento]{AB}$.

% doc: «lado A B»
69. Con concept true, el valor por defecto: $\spside[concept=true]{AB}$.

% doc: «A B»
70. Con concept false: $\spside[concept=false]{AB}$.

% doc: «A B»
71. Con concept false y trazo: $\spside[concept=false,read-sty=trazo]{AB}$.

% doc: «hipotenusa A B»
72. Con read-arg: $\spside[read-arg={hipotenusa}]{AB}$.

% doc: «lado opuesto A B»
73. Con read-arg de varias palabras: $\spside[read-arg={lado opuesto}]{AB}$.

% doc: «la hipotenusa A B»
74. Con read-arg y prefix: $\spside[read-arg={hipotenusa},prefix={la}]{AB}$.

% doc: «hipotenusa A B»
75. Con read-arg y concept false: $\spside[read-arg={hipotenusa},concept=false]{AB}$.

% doc: «lado A B»
76. Con silent false, el valor por defecto: $\spside[silent=false]{AB}$.

% doc: «»
77. Con silent true: $\spside[silent=true]{AB}$.

% doc: «»
78. Con silent true y subíndices: $\spside[silent=true]{A_{1}B_{2}}$.

% doc: «el segmento a subíndice 1 b subíndice 2»
79. Con todas las llaves: $\spside[read-cmd=mathcat,read-sty=segmento,prefix={el},concept=true,silent=false]{A_{1}B_{2}}$.

\section*{E. Comando spmside}

% doc: «medida del lado A B»
80. Dos puntos juntos: $\spmside{AB}$.

% doc: «medida del lado A B»
81. Dos puntos con coma y espacio: $\spmside{A, B}$.

% doc: «medida del lado A B»
82. Dos puntos con coma: $\spmside{A,B}$.

% doc: «medida del lado A prima B prima»
83. Con primas: $\spmside{A'B'}$.

% doc: «medida del lado A sub uno B sub dos»
84. Con subíndices: $\spmside{A_{1}B_{2}}$.

% doc: «medida del lado A sub uno B sub dos»
85. Con subíndices y coma: $\spmside{A_{1},B_{2}}$.

% doc: «medida del lado B C»
86. Con el vértice: $\spmside{BC}$.

% doc: «medida del lado A B»
87. Con read-cmd school, el valor por defecto: $\spmside[read-cmd=school]{AB}$.

% doc: «medida del lado a b»
88. Con read-cmd mathcat: $\spmside[read-cmd=mathcat]{AB}$.

% doc: «medida del lado A B»
89. Con read-sty lado: $\spmside[read-sty=lado]{AB}$.

% doc: «medida del segmento A B»
90. Con read-sty segmento: $\spmside[read-sty=segmento]{AB}$.

% doc: «medida del trazo A B»
91. Con read-sty trazo: $\spmside[read-sty=trazo]{AB}$.

% doc: «el medida del lado A B»
92. Con prefix: $\spmside[prefix={el}]{AB}$.

% doc: «el medida del segmento A B»
93. Con prefix y read-sty segmento: $\spmside[prefix={el},read-sty=segmento]{AB}$.

% doc: «medida del lado A B»
94. Con concept true, el valor por defecto: $\spmside[concept=true]{AB}$.

% doc: «A B»
95. Con concept false: $\spmside[concept=false]{AB}$.

% doc: «A B»
96. Con concept false y trazo: $\spmside[concept=false,read-sty=trazo]{AB}$.

% doc: «medida de la cuerda A B»
97. Con read-arg y concept false: $\spmside[concept=false,read-arg={medida de la cuerda}]{AB}$.

% doc: «medida del hipotenusa A B»
98. Con read-arg y concept true: $\spmside[read-arg={hipotenusa}]{AB}$.

% doc: «la medida de la cuerda A B»
99. Con read-arg y prefix: $\spmside[concept=false,read-arg={medida de la cuerda},prefix={la}]{AB}$.

% doc: «medida del lado A B»
100. Con print-m false, el valor por defecto: $\spmside[print-m=false]{AB}$.

% doc: «medida del lado A B»
101. Con print-m true: $\spmside[print-m=true]{AB}$.

% doc: «medida del lado A sub uno B sub dos»
102. Con print-m true y subíndices: $\spmside[print-m=true]{A_{1}B_{2}}$.

% doc: «medida del lado A B»
103. Con silent false, el valor por defecto: $\spmside[silent=false]{AB}$.

% doc: «»
104. Con silent true: $\spmside[silent=true]{AB}$.

% doc: «»
105. Con silent true y subíndices: $\spmside[silent=true]{A_{1}B_{2}}$.

% doc: «el medida del segmento a subíndice 1 b subíndice 2»
106. Con todas las llaves: $\spmside[read-cmd=mathcat,read-sty=segmento,prefix={el},concept=true,silent=false,print-m=true]{A_{1}B_{2}}$.

\section*{F. En fórmulas}

% doc: «recta A B es paralelo a recta C D»
107. Rectas paralelas: $\sprecta{AB} \parallel \sprecta{CD}$.

% doc: «recta m es perpendicular a recta n»
108. Rectas perpendiculares: $\sprecta{m} \perp \sprecta{n}$.

% doc: «rayo O A unión rayo O B»
109. Dos rayos con origen común: $\sprayo{OA} \cup \sprayo{OB}$.

% doc: «lado A B es igual a lado C D»
110. Segmentos iguales: $\spside{AB} = \spside{CD}$.

% doc: «medida del lado A B es igual a medida del lado C D»
111. Medidas iguales: $\spmside{AB} = \spmside{CD}$.

% doc: «medida del lado A B más medida del lado B C es mayor que medida del lado A C»
112. Desigualdad triangular: $\spmside{AB} + \spmside{BC} > \spmside{AC}$.

% doc: «medida del lado A B es igual a medida del lado C D»
113. Medida con m impresa: $\spmside[print-m=true]{AB} = \spmside[print-m=true]{CD}$.

% doc: «medida del lado A B más medida del lado B C es igual a medida del lado A C»
114. En fórmula destacada:
\[ \spmside{AB} + \spmside{BC} = \spmside{AC} \].

\end{document}
